\documentclass[12pt]{article}
\usepackage[a4paper,margin=1in]{geometry}
\usepackage{amsmath,amssymb}
\usepackage{enumitem}
\usepackage{mathtools}
\usepackage{hyperref}
\usepackage{titlesec}
\usepackage{amsthm}

\title{Topology MATM36: Course Summary of Important Results}
\author{Simon Gustafsson}
\date{}

\newtheorem{theorem}{Theorem}
\newtheorem{corollary}{Corollary}
\newtheorem{definition}{Definition}
\newtheorem{remark}{Remark}


\begin{document}
\maketitle

\section*{Metric Spaces}

\begin{definition}
Let $X$ be a topological space and $A \subseteq X$. A point $x \in X$ is called an \emph{adherent point} of $A$ if every open neighbourhood $U$ of $x$ satisfies
\[
U \cap A \neq \varnothing.
\]
\end{definition}

In the special case where $X$ is a metric space with metric $d$, this is equivalent to requiring that for every $r>0$,
\[
B_r(x) \cap A \neq \varnothing.
\]

\begin{definition}
Let $X$ be a topological space and $A \subseteq X$. The set $A$ is called \emph{dense} in $X$ if every nonempty open set $U \subseteq X$ satisfies
\[
U \cap A \neq \varnothing.
\]
Equivalently, $A$ is dense in $X$ if $\overline{A} = X$.
\end{definition}

\begin{definition}
Let $X$ be a topological space and $A \subseteq X$. 
The set $A$ is called \emph{nowhere dense} in $X$ if
\[
\operatorname{int}(\overline{A}) = \varnothing.
\]
\end{definition}

\subsection*{Baire Category Theorem}

\begin{theorem}[Baire Category Theorem]
Let $(X,d)$ be a complete metric space and let $\{U_n\}_{n=1}^{\infty}$ be a sequence of dense open subsets of $X$. Then \(\bigcap_{n=1}^{\infty} U_n\) is dense in $X$.
\end{theorem}

\begin{proof}
Fix a nonempty open ball $B(x,\varepsilon)\subset X$.

Since $U_1$ is dense and open, choose $y_1\in U_1\cap B(x,\varepsilon)$ and $r_1>0$ such that
\[
\overline{B}(y_1,r_1)\subset U_1\cap B(x,\varepsilon),\qquad r_1<1.
\]
Inductively, having chosen $y_n,r_n$, use that $U_{n+1}$ is dense and open to pick
$y_{n+1}\in U_{n+1}\cap B(y_n,r_n)$ and $r_{n+1}>0$ such that
\[
\overline{B}(y_{n+1},r_{n+1})\subset U_{n+1}\cap B(y_n,r_n),\qquad r_{n+1}<\frac1{n+1}.
\]
Then the closed balls are nested:
\[
\overline{B}(y_{n+1},r_{n+1})\subset \overline{B}(y_n,r_n)\subset \cdots \subset \overline{B}(y_1,r_1)\subset B(x,\varepsilon),
\]
and for $m>n$ we have $y_m\in B(y_n,r_n)$, hence $d(y_m,y_n)<r_n\to 0$. Thus $(y_n)$ is Cauchy, so by completeness $y_n\to y\in X$.

For each fixed $n$, all tails of the sequence lie in $\overline{B}(y_n,r_n)$, so the limit satisfies
$y\in \overline{B}(y_n,r_n)\subset U_n$. Hence $y\in \bigcap_{n=1}^\infty U_n$, and also
$y\in \overline{B}(y_1,r_1)\subset B(x,\varepsilon)$. Therefore
\[
B(x,\varepsilon)\cap\bigcap_{n=1}^\infty U_n\neq\varnothing,
\]
so the intersection is dense.
\end{proof}


\begin{remark}
From the definitions it follows that a set $B \subseteq X$ satisfies 
\(\operatorname{int}(B)=\varnothing \Longleftrightarrow X\setminus B \text{ is dense in } X.\)
In particular, if $A$ is nowhere dense, then
\( X\setminus \overline{A}\) 
is open and dense in $X$.
\end{remark}

\begin{corollary}
Let $(X,d)$ be a complete metric space and let $\{F_n\}_{n=1}^{\infty}$ be a sequence of nowhere dense subsets of $X$. Then
\[
\operatorname{int}\!\left(\bigcup_{n=1}^{\infty} F_n\right)=\varnothing.
\]
\end{corollary}

\begin{proof}
Since $F_n$ is nowhere dense, $U_n:=X\setminus\overline{F_n}$ is open and dense; by Baire, $\bigcap_{n=1}^\infty U_n = \bigcup_{n=1}^\infty \overline{F_n}$ is dense, hence nonempty in every nonempty open set.
\end{proof}

\subsection*{Finite Products of Metric Spaces}

\begin{definition}
Let $(X_k,d_k)$, $k=1,\dots,n$, be metric spaces and set
\[
X = X_1 \times \cdots \times X_n.
\]
A metric $d$ on $X$ is said to be a \emph{product metric} if convergence in $(X,d)$ is coordinatewise,
\[
x^{(m)} = (x_1^{(m)},\dots,x_n^{(m)}) \to x=(x_1,\dots,x_n)
\]
in $X$ if and only if for each $k=1,\dots,n$,
\[
x_k^{(m)} \to x_k \quad \text{in } X_k.
\]
\end{definition}

\begin{theorem}
Suppose $d$ is a metric on $X=X_1\times\cdots\times X_n$ satisfying the above coordinatewise convergence property. Then the open sets in $(X,d)$ are precisely the unions of product sets
\[
U_1 \times \cdots \times U_n,
\]
where each $U_k$ is open in $X_k$.
\end{theorem}

\begin{proof}
First, if $E_k \subset X_k$ is closed for each $k$, then
\[
E_1 \times \cdots \times E_n
\]
is closed in $X$. Indeed, if $x^{(m)} \in E_1 \times \cdots \times E_n$ and $x^{(m)} \to x$ in $X$, then by coordinatewise convergence,
$x_k^{(m)} \to x_k$ in $X_k$. Since each $E_k$ is closed, $x_k \in E_k$, hence
$x \in E_1 \times \cdots \times E_n$.

Now let $U_k \subset X_k$ be open. Then
\[
V^{(k)} := X_1 \times \cdots \times X_{k-1} \times U_k \times X_{k+1} \times \cdots \times X_n
\]
is open in $X$, since its complement is of the form
\[
X_1 \times \cdots \times (X_k \setminus U_k) \times \cdots \times X_n,
\]
which is closed by the previous remark. Hence
\[
U_1 \times \cdots \times U_n
=
\bigcap_{k=1}^n V^{(k)}
\]
is open in $X$, and any union of such products is open.

Conversely, let $U \subset X$ be open and $x=(x_1,\dots,x_n)\in U$.
We show there exist open sets $U_k \subset X_k$ such that
\[
x \in U_1 \times \cdots \times U_n \subset U.
\]
Suppose not. Then for each $m\in\mathbb{N}$,
\[
B_{X_1}(x_1,1/m) \times \cdots \times B_{X_n}(x_n,1/m)
\not\subset U.
\]
Thus there exists $x^{(m)} \notin U$ with
\[
d_k(x_k^{(m)},x_k) < \frac1m, \qquad k=1,\dots,n.
\]
Then $x_k^{(m)} \to x_k$ for each $k$, hence $x^{(m)} \to x$ in $X$.
Since $X\setminus U$ is closed and each $x^{(m)} \in X\setminus U$,
we conclude $x \in X\setminus U$, a contradiction.
Therefore such open sets $U_k$ exist, and $U$ is a union of product sets.
\end{proof}

\end{document}